\begin{example}
    e.g. Encoding a logic puzzle. \vspace{7pt}

    We consider the following clues:
    \begin{enumerate}
        \item Good-natured tenured professors are dynamic.
        \item Grumpy student advisors play slot machines.
        \item Smokers wearing a cap are phlegmatic.
        \item Comical student advisors are professors.
        \item Smoking untenured members are nervous.
        \item Phlegmatic tenured members wearing caps are comical.
        \item Student advisors who are not stock market players are scholars.
        \item Relaxed student advisors are creative.
        \item Creative scholars who do not play slot machines wear caps.
        \item Nervous smokers play slot machines.
        \item Student advisors who play slot machines do not smoke.
        \item Creative good-natured stock market players wear caps.
    \end{enumerate} \vspace{7pt}

    Can we conclude that \textbf{no student advisor is smoking}? \vspace{7pt}

    In order to answer the previous question, we must give a name for every notion to be formalized. \vspace{7pt}

    \begin{center}
        \begin{tabular}{|c|c|c|}
            \hline
            \textbf{Name} & \textbf{Meaning} & \textbf{Opposite} \\ \hline
            $a$ & good-natured & grumpy \\ \hline
            $b$ & tenured & \\ \hline
            $c$ & professor & \\ \hline
            $d$ & dynamic & phlegmatic \\ \hline
            $e$ & wearing a cap & \\ \hline
            $f$ & smoke & \\ \hline
            $g$ & comical & \\ \hline
            $h$ & relaxed & nervous \\ \hline
            $i$ & play stock market & \\ \hline
            $j$ & scholar & \\ \hline
            $k$ & creative & \\ \hline
            $l$ & play slot machine & \\ \hline
            $m$ & student advisor & \\ \hline
        \end{tabular}
    \end{center} \vspace{7pt}

    Defined a name for every notion, now we are able to express natural language clues in propositional formulas. For instance, the first clue:
    $$ \text{Good-natured tenured professors are dynamic} $$ 
    becomes:
    $$ (a \wedge b \wedge c) \rightarrow d $$
    where:
    \begin{itemize}
        \renewcommand{\labelitemi}{-}
        \item $a \rightarrow $ good-natured
        \item $b \rightarrow $ tenured
        \item $c \rightarrow$ professors
        \item $d \rightarrow$ dynamic
    \end{itemize} \vspace{7pt}

    Therefore, the whole set of clues becomes as follows:
    \begin{enumerate}
        \item $(a \wedge b \wedge c) \to d$
        \item $(\neg a \wedge m) \to l$
        \item $(f \wedge e) \to \neg d$
        \item $(g \wedge m) \to c$
        \item $(f \wedge \neg b) \to \neg h$
        \item $(\neg d \wedge b \wedge e) \to g$
        \item $(\neg i \wedge m) \to j$
        \item $(h \wedge m) \to k$
        \item $(k \wedge j \wedge \neg l) \to e$
        \item $(\neg h \wedge f) \to l$
        \item $(l \wedge m) \to \neg f$
        \item $(k \wedge a \wedge i) \to e$
    \end{enumerate} \vspace{7pt}

    After this process of encoding, the original question now is:
    $$ \neg (m \wedge f)? $$
    Giving to the SAT solver the previous question and the set of propositional formulas listed before, it proves that the formula $\neg (m \wedge f)$ is UNSAT,
    so we conclude saying that there is no student advisor that is smoking given this set of clues.
\end{example}